\documentclass[../../../include/open-logic-section]{subfiles}

\begin{document}

\olfileid{sth}{card-arithmetic}{opps}
\olsection{تعريف العمليات الأساسية}

لما كنا غير محتاجين إلى حفظ الترتيب، كان تعريف حساب الكاردينالات
أيسر من تعريف حساب الأعداد الترتيبية بعض اليسر. ونعرف الجمع والضرب
والأس جميعًا.

\begin{defn}
إذا كان~$\cardfont{\OLId{latin-ordinary}{a}}$ و$\cardfont{\OLId{latin-ordinary}{b}}$ كاردينالين، فنعرف:
\begin{align*}
	\cardfont{\OLId{latin-ordinary}{a}} \cardplus \cardfont{\OLId{latin-ordinary}{b}} &\defis
	\card{\cardfont{\OLId{latin-ordinary}{a}} \disjointsum \cardfont{\OLId{latin-ordinary}{b}}}\\
	\cardfont{\OLId{latin-ordinary}{a}} \cardtimes \cardfont{\OLId{latin-ordinary}{b}} &\defis
	\card{\cardfont{\OLId{latin-ordinary}{a}} \times \cardfont{\OLId{latin-ordinary}{b}}}\\
	\cardexpo{\cardfont{\OLId{latin-ordinary}{a}}}{\cardfont{\OLId{latin-ordinary}{b}}} &\defis
	\card{\funfromto{\cardfont{\OLId{latin-ordinary}{b}}}{\cardfont{\OLId{latin-ordinary}{a}}}}
\end{align*}
حيث
$\funfromto{\OLId{latin-looped}{X}}{\OLId{latin-looped}{Y}} = \Setabs{\OLId{latin-ordinary}{f}}{\OLId{latin-ordinary}{f} \text{ دالة } \OLId{latin-looped}{X} \to \OLId{latin-looped}{Y}}$.
(ووجود~$\funfromto{\OLId{latin-looped}{X}}{\OLId{latin-looped}{Y}}$ لأي مجموعتين~$\OLId{latin-looped}{X}$ و$\OLId{latin-looped}{Y}$ سهل الإثبات، وهو
متروك تمرينًا.)
\end{defn}

\begin{prob}
أثبت في~$\Zminus$ أن~$\funfromto{\OLId{latin-looped}{X}}{\OLId{latin-looped}{Y}}$ موجودة لأي مجموعتين~$\OLId{latin-looped}{X}$
و~$\OLId{latin-looped}{Y}$. وبالعمل في~$\ZF$، احسب~$\setrank{\funfromto{\OLId{latin-looped}{X}}{\OLId{latin-looped}{Y}}}$ من
$\setrank{\OLId{latin-looped}{X}}$ و$\setrank{\OLId{latin-looped}{Y}}$، على مثال
\olref[sth][ord-arithmetic][using-addition]{rankcomputation}.
\end{prob}

ولنشرح وجه التعريف. فأما الجمع فيستعمل الجمع المنفصل،~$\disjointsum$،
على تعريف~\olref[ord-arithmetic][add]{defdissum}؛ ومن اليسير رؤية
موافقته للحساب في الحالات المتناهية. وأما الضرب، فقد أثبت
\olref[sfr][set][pai]{cardnmprod} أن المجموعة~$\OLId{latin-looped}{A}$ إذا كان لها~$\OLId{latin-ordinary}{n}$
عناصر، والمجموعة~$\OLId{latin-looped}{B}$ إذا كان لها~$\OLId{latin-ordinary}{m}$ عناصر، كان لـ$\OLId{latin-looped}{A} \times \OLId{latin-looped}{B}$
عدد~$\OLId{latin-ordinary}{n} \cdot \OLId{latin-ordinary}{m}$ من العناصر. فتعريفنا يعمم هذا المعنى على ما فوق
المتناهي. والأس على القياس نفسه، ينقل الفكرة من المتناهي إلى ما فوقه.
بل إن حساب الكاردينالات ما فوق المتناهي أشبه بالحساب المعتاد من بعض
الوجوه من حساب الأعداد الترتيبية ما فوق المتناهي.

\begin{prop}\ollabel{cardplustimescommute}
العمليتان~$\cardplus$ و$\cardtimes$ إبداليتان وتجميعيتان.
\end{prop}

\begin{proof}
يكفي في الإبدالية، بموجب
\olref[cardinals][cardsasords]{lem:CardinalsBehaveRight}، أن نلاحظ
$\cardeq{(\cardfont{\OLId{latin-ordinary}{a}} \disjointsum
\cardfont{\OLId{latin-ordinary}{b}})}{(\cardfont{\OLId{latin-ordinary}{b}} \disjointsum \cardfont{\OLId{latin-ordinary}{a}})}$ وأن
$\cardeq{(\cardfont{\OLId{latin-ordinary}{a}} \times \cardfont{\OLId{latin-ordinary}{b}})}{(\cardfont{\OLId{latin-ordinary}{b}} \times
\cardfont{\OLId{latin-ordinary}{a}})}$. وأما التجميعية فنتركها تمرينًا.
\end{proof}

\begin{prob}
أثبت أن~$\cardplus$ و$\cardtimes$ تجميعيتان.
\end{prob}

\begin{prop}
$\OLId{latin-looped}{A}$ غير متناهية إذا وفقط إذا كان
$\card{\OLId{latin-looped}{A}} \cardplus \OLMathNumeral{1} = \OLMathNumeral{1} \cardplus \card{\OLId{latin-looped}{A}} = \card{\OLId{latin-looped}{A}}$.
\end{prop}

\begin{proof}
كما في~\olref[cardinals][classing]{generalinfinitycharacter}، تتبع
النتيجة من~\olref[ord-arithmetic][using-addition]{ordinfinitycharacter}
و\olref[cardinals][cardsasords]{lem:CardinalsBehaveRight}.
\end{proof}

ومن هنا احتيج إلى رموز تفصل جمع الأعداد الترتيبية وضربها عن جمع
الكاردينالات وضربها؛ فهما عمليتان \emph{مختلفتان} حقًا. والنتيجتان
الآتيتان تبينان اختلاف الأسين كذلك. وتذكر أن
\olref[z][infinity-again]{defnomega} يقتضي~$\OLMathNumeral{2} = \{\OLMathNumeral{0}, \OLMathNumeral{1}\}$.

\begin{lem}\ollabel{lem:SizePowerset2Exp}
لكل~$\OLId{latin-looped}{A}$، $\card{\Pow{\OLId{latin-looped}{A}}} = \cardexpo{\OLMathNumeral{2}}{\card{\OLId{latin-looped}{A}}}$.
\end{lem}

\begin{proof}
لكل مجموعة جزئية~$\OLId{latin-looped}{B} \subseteq \OLId{latin-looped}{A}$، لتكن
$\chi_\OLId{latin-looped}{B} \in \funfromto{\OLId{latin-looped}{A}}{\OLMathNumeral{2}}$ معطاة هكذا:
\begin{align*}
	\chi_{\OLId{latin-looped}{B}}(\OLId{latin-ordinary}{x}) &\defis
	\begin{cases}
		\OLMathNumeral{1} & \text{إذا كان }\OLId{latin-ordinary}{x}\in \OLId{latin-looped}{B}\\
		\OLMathNumeral{0} & \text{فيما عدا ذلك.}
	\end{cases}
\end{align*}
ثم ضع~$\OLId{latin-ordinary}{f}(\OLId{latin-looped}{B}) = \chi_\OLId{latin-looped}{B}$؛ فيعرف ذلك !!a{bijection}
$\OLId{latin-ordinary}{f} \colon \Pow{\OLId{latin-looped}{A}} \to \funfromto{\OLId{latin-looped}{A}}{\OLMathNumeral{2}}$. ومن ثم
$\cardeq{\Pow{\OLId{latin-looped}{A}}}{\funfromto{\OLId{latin-looped}{A}}{\OLMathNumeral{2}}}$، ويتبعه
$\cardeq{\Pow{\OLId{latin-looped}{A}}}{\funfromto{\card{\OLId{latin-looped}{A}}}{\OLMathNumeral{2}}}$. ولذلك
$\card{\Pow{\OLId{latin-looped}{A}}} = \card{\funfromto{\card{\OLId{latin-looped}{A}}}{\OLMathNumeral{2}}} =
\cardexpo{\OLMathNumeral{2}}{\card{\OLId{latin-looped}{A}}}$.
\end{proof}

وهذا البرهان الوجيز يجمع في باطنه بحث
\olref[sfr][siz][red-alt]{sec}. فقد رددنا هنالك عدم قابلية
$\Pow{\omega}$ للعد إلى عدم قابلية مجموعة السلاسل الثنائية غير
المتناهية~$\Bin^\omega$. وليست~$\Bin^{\omega}$ إلا
$\funfromto{\omega}{\OLMathNumeral{2}}$. وقد ظهر الآن أن مسلك الاستدلال في
\olref[sfr][siz][red-alt]{sec} نفسه يصح إذا عوضنا أي مجموعة~$\OLId{latin-looped}{A}$
عن~$\omega$. وتفيد النتيجة أيضًا حقيقة قريبة عن أس
الكاردينالات.

\begin{cor}\ollabel{cantorcor}
لكل كاردينال~$\cardfont{\OLId{latin-ordinary}{a}}$،
$\cardfont{\OLId{latin-ordinary}{a}} < \cardexpo{\OLMathNumeral{2}}{\cardfont{\OLId{latin-ordinary}{a}}}$.
\end{cor}

\begin{proof}
ذلك من مبرهنة كانتور~(\olref[sfr][siz][car]{thm:cantor}) ومن
\olref{lem:SizePowerset2Exp}.
\end{proof}
\noindent
إذن~$\omega < \cardexpo{\OLMathNumeral{2}}{\omega}$. وهذا حكم في أس
\emph{الكاردينالات}؛ فقابله بأس \emph{الأعداد الترتيبية}، حيث
$\omega = \ordexpo{\OLMathNumeral{2}}{\omega}$؛ انظر
\olref[ord-arithmetic][expo]{sec}.

وبمناسبة أس الكاردينالات يمكن أن نضبط أكثر الوجه الذي تكون به
$\Real$ !!{nonenumerable}.

\begin{thm}\ollabel{continuumis2aleph0}
$\card{\Real} = \cardexpo{\OLMathNumeral{2}}{\omega}$.
\end{thm}

\begin{proof}[مخطط البرهان]
لهذا براهين كثيرة. وأقربها أن نثبت
$\cardle{\Pow{\omega}}{\Real}$ و$\cardle{\Real}{\Pow{\omega}}$، ثم
نستعمل شرودر--برنشتاين فنحصل على
$\cardeq{\Real}{\Pow{\omega}}$، ثم
\olref[card-arithmetic][opps]{lem:SizePowerset2Exp} فنستنتج
$\card{\Real} = \cardexpo{\OLMathNumeral{2}}{\omega}$. ويبقى تمرين نافع: عرّف
دالتين حقنيتين~$\OLId{latin-ordinary}{f} \colon \Pow{\omega} \to \Real$ و
$\OLId{latin-ordinary}{g} \colon \Real \to \Pow{\omega}$.
\end{proof}

\begin{prob}
أكمل برهان~\olref[sth][card-arithmetic][opps]{continuumis2aleph0} بإظهار
أن~$\cardle{\Pow{\omega}}{\Real}$ و
$\cardle{\Real}{\Pow{\omega}}$.
\end{prob}

\end{document}
